\chapter{Módulo 5 --- Transformar datos con dplyr}

\begin{itemize}
\tightlist
\item
  \textbf{Duración estimada:} 1.5 h teoría + 3.5 h práctica
\item
  \textbf{Objetivos de aprendizaje.} Al finalizar, la persona podrá:

  \begin{enumerate}
  \def\labelenumi{\arabic{enumi}.}
  \tightlist
  \item
    \textbf{Aplicar} los verbos de \passthrough{\lstinline!dplyr!}
    (\passthrough{\lstinline!filter!}, \passthrough{\lstinline!select!},
    \passthrough{\lstinline!mutate!}, \passthrough{\lstinline!arrange!},
    \passthrough{\lstinline!summarise!}) encadenados con
    \passthrough{\lstinline!|>!}.
  \item
    \textbf{Calcular} indicadores agregados por grupo con
    \passthrough{\lstinline!group\_by()!} /
    \passthrough{\lstinline!.by!}.
  \item
    \textbf{Combinar} tablas con uniones
    (\passthrough{\lstinline!left\_join!},
    \passthrough{\lstinline!inner\_join!},
    \passthrough{\lstinline!anti\_join!}).
  \item
    \textbf{Usar} \passthrough{\lstinline!lubridate!} y
    \passthrough{\lstinline!stringr!} para trabajar con fechas y texto.
  \item
    \textbf{Analizar} el desempeño por canal, región y producto a partir
    de datos crudos.
  \end{enumerate}
\end{itemize}

\section{Contenidos}

\begin{enumerate}
\def\labelenumi{\arabic{enumi}.}
\tightlist
\item
  Los cinco verbos principales y el \emph{pipe}.
\item
  \passthrough{\lstinline!mutate()!} con
  \passthrough{\lstinline!if\_else()!} y
  \passthrough{\lstinline!case\_when()!}.
\item
  \passthrough{\lstinline!group\_by()!} +
  \passthrough{\lstinline!summarise()!};
  \passthrough{\lstinline!count()!}; agrupación temporal con
  \passthrough{\lstinline!.by!}.
\item
  Funciones de ventana: \passthrough{\lstinline!lag()!},
  \passthrough{\lstinline!cumsum()!}, \passthrough{\lstinline!rank()!},
  \passthrough{\lstinline!slice\_max()!}.
\item
  Uniones de tablas y llaves.
\item
  Fechas con \passthrough{\lstinline!lubridate!}:
  \passthrough{\lstinline!year()!}, \passthrough{\lstinline!month()!},
  \passthrough{\lstinline!floor\_date()!}.
\item
  Texto con \passthrough{\lstinline!stringr!}:
  \passthrough{\lstinline!str\_detect()!},
  \passthrough{\lstinline!str\_to\_upper()!},
  \passthrough{\lstinline!str\_replace()!}.
\end{enumerate}

\section{Desarrollo}

\begin{lstlisting}[language=R]
library(tidyverse)
ventas  <- read_csv("datos/ventas.csv", show_col_types = FALSE)
tiendas <- read_csv("datos/tiendas.csv", show_col_types = FALSE)
\end{lstlisting}

\subsection{Los verbos básicos}

\begin{lstlisting}[language=R]
ventas |>
  filter(producto == "Tablet", unidades >= 8) |>
  select(fecha, tienda_id, unidades) |>
  arrange(desc(unidades))

ventas <- ventas |>
  mutate(ingreso = precio * unidades)
\end{lstlisting}

\subsection{Resúmenes por grupo}

\begin{lstlisting}[language=R]
ventas |>
  group_by(producto) |>
  summarise(
    unidades = sum(unidades, na.rm = TRUE),
    ingreso  = sum(ingreso,  na.rm = TRUE),
    ticket_promedio = ingreso / unidades
  ) |>
  arrange(desc(ingreso))

ventas |> count(producto)                         # conteo rápido
\end{lstlisting}

\subsection{Uniones}

\begin{lstlisting}[language=R]
ventas_det <- ventas |>
  left_join(tiendas, by = "tienda_id")

ventas_det |>
  summarise(ingreso = sum(ingreso, na.rm = TRUE), .by = c(canal, region)) |>
  arrange(canal, desc(ingreso))

# ¿Hay ventas de tiendas que no existen en el catálogo?
ventas |> anti_join(tiendas, by = "tienda_id") |> nrow()   # 0
\end{lstlisting}

\subsection{Fechas}

\begin{lstlisting}[language=R]
mensual <- ventas_det |>
  mutate(mes = floor_date(fecha, "month")) |>
  summarise(ingreso = sum(ingreso, na.rm = TRUE), .by = c(mes, canal)) |>
  arrange(canal, mes) |>
  group_by(canal) |>
  mutate(var_pct = (ingreso / lag(ingreso) - 1) * 100) |>
  ungroup()
mensual
\end{lstlisting}

\subsection{Texto}

\begin{lstlisting}[language=R]
ventas_det |>
  filter(str_detect(producto, "Smartphone")) |>
  mutate(producto = str_to_upper(producto)) |>
  distinct(producto)
\end{lstlisting}

\section{Actividades y ejercicios}

\begin{enumerate}
\def\labelenumi{\arabic{enumi}.}
\tightlist
\item
  ¿Cuál fue la \textbf{tienda} con mayor ingreso anual? Muestra el top 3
  con su canal y región.
\item
  Calcula la \textbf{participación (\%)} de cada canal en el ingreso
  total.
\item
  Para cada producto, encuentra la \textbf{semana} de mayores unidades
  vendidas (\passthrough{\lstinline!slice\_max()!}).
\item
  Compara el ingreso promedio semanal de \textbf{noviembre--diciembre}
  vs.~el resto del año por producto.
\item
  Reto: con \passthrough{\lstinline!anti\_join()!}, verifica que cada
  combinación tienda--semana tenga los 5 productos.
\end{enumerate}

\subsection{Soluciones}

\begin{lstlisting}[language=R]
# 1
ventas_det |>
  summarise(ingreso = sum(ingreso, na.rm = TRUE), .by = c(tienda_id, canal, region)) |>
  slice_max(ingreso, n = 3)

# 2
ventas_det |>
  summarise(ingreso = sum(ingreso, na.rm = TRUE), .by = canal) |>
  mutate(participacion = ingreso / sum(ingreso) * 100)

# 3
ventas |>
  summarise(unidades = sum(unidades, na.rm = TRUE), .by = c(producto, fecha)) |>
  group_by(producto) |>
  slice_max(unidades, n = 1, with_ties = FALSE)

# 4
ventas |>
  mutate(temporada = if_else(month(fecha) %in% 11:12, "Nov-Dic", "Resto")) |>
  summarise(ingreso_sem = sum(ingreso, na.rm = TRUE), .by = c(producto, temporada, fecha)) |>
  summarise(promedio = mean(ingreso_sem), .by = c(producto, temporada)) |>
  pivot_wider(names_from = temporada, values_from = promedio) |>
  mutate(incremento_pct = (`Nov-Dic` / Resto - 1) * 100)

# 5
esperado <- expand_grid(
  tienda_id = unique(ventas$tienda_id),
  fecha     = unique(ventas$fecha),
  producto  = unique(ventas$producto)
)
esperado |> anti_join(ventas, by = c("tienda_id", "fecha", "producto")) |> nrow()  # 0
\end{lstlisting}

\section{Evaluación (sumativa parcial)}

\textbf{Mini-reporte de KPIs:} script que produzca una tabla con
ingreso, unidades y ticket promedio por canal y región, más la variación
mensual por canal. Se evalúa: resultados correctos (40 \%), uso
idiomático de \passthrough{\lstinline!dplyr!} (30 \%), legibilidad (30
\%).

\section{Referencias verificadas}

\begin{itemize}
\tightlist
\item
  \textbf{Libro:} Wickham, H., Çetinkaya-Rundel, M., \& Grolemund, G.
  (2023). \emph{R for Data Science} (2.ª ed.). O'Reilly Media. ISBN
  978-1-4920-9740-2. Español: \url{https://davidrsch.github.io/r4ds-es/}
  {[}ES{]} {[}OA{]} --- caps. ``Transformación de datos'', ``Uniones'',
  ``Fechas y horas'', ``Cadenas''.
\item
  \textbf{Web:} R para Ciencia de Datos (traducción de la 1.ª ed.~por la
  comunidad latinoamericana de R). \url{https://es.r4ds.hadley.nz/}
  {[}ES{]} {[}OA{]} --- alternativa en español.
\item
  \textbf{Artículo:} Wickham, H., et al.~(2019). Welcome to the
  tidyverse. \emph{Journal of Open Source Software, 4}(43), 1686.
  \url{https://doi.org/10.21105/joss.01686} {[}EN{]} {[}OA{]}
\item
  \textbf{Video:} Gomila Salas, J. G. (Frogames). \emph{Curso completo
  de R para Data Science con Tidyverse} {[}Video{]}. YouTube.
  \url{https://www.youtube.com/watch?v=9NJyhs5PlGQ} {[}ES{]}
  \emph{{[}canal y fecha POR VERIFICAR{]}}
\end{itemize}
